\section*{NeurIPS Paper Checklist}

\begin{enumerate}

\item {\bf Claims}
    \item[] Question: Do the main claims made in the abstract and introduction accurately reflect the paper's contributions and scope?
    \item[] Answer: \answerYes{}
    \item[] Justification: The abstract and introduction state the forward/control-pass, inverse-fail dissociation, the ours-equals-oracle match, the gap that widens with horizon, and the cross-domain result; all are supported by Sections~\ref{sec:results}--\ref{sec:wall}, and scope and caveats are in Section~\ref{sec:limitations}.

\item {\bf Limitations}
    \item[] Question: Does the paper discuss the limitations of the work performed by the authors?
    \item[] Answer: \answerYes{}
    \item[] Justification: Section~\ref{sec:limitations} discusses the abstract toy domains, exact-equality scoring, the measured enumeration wall (MAP feasible only to $L\approx 12$ / $8\times8$), the small LLM sample sizes, and that o3 is not weak in general.

\item {\bf Theory assumptions and proofs}
    \item[] Question: For each theoretical result, does the paper provide the full set of assumptions and a complete (and correct) proof?
    \item[] Answer: \answerNA{}
    \item[] Justification: The paper is empirical; it states the MAP estimator and its likelihood (Section~\ref{sec:setup}) but contains no formal theorems requiring proof.

\item {\bf Experimental result reproducibility}
    \item[] Question: Does the paper fully disclose all the information needed to reproduce the main experimental results of the paper to the extent that it affects the main claims and/or conclusions of the paper (regardless of whether the code and data are provided or not)?
    \item[] Answer: \answerYes{}
    \item[] Justification: Section~\ref{sec:setup} specifies the model, training, candidate sets, noise channel, oracle, and metrics; Appendix~\ref{app:prompts} gives the LLM prompt and extraction; code and cached transcripts are released with a one-command \texttt{reproduce.sh}.

\item {\bf Open access to data and code}
    \item[] Question: Does the paper provide open access to the data and code, with sufficient instructions to faithfully reproduce the main experimental results, as described in supplemental material?
    \item[] Answer: \answerYes{}
    \item[] Justification: The code, instance generators, and the cached o3/Sonnet transcripts are released; \texttt{reproduce.sh} regenerates every CPU result and aggregates the cached LLM numbers without any API access.

\item {\bf Experimental setting/details}
    \item[] Question: Does the paper specify all the training and test details (e.g., data splits, hyperparameters, how they were chosen, type of optimizer) necessary to understand the results?
    \item[] Answer: \answerYes{}
    \item[] Justification: Section~\ref{sec:setup} gives the MLP architecture, optimizer, learning rate, training steps, training widths, sample sizes, seeds, noise level, and horizons; the LLM decoding settings are stated and the prompts are in Appendix~\ref{app:prompts}.

\item {\bf Experiment statistical significance}
    \item[] Question: Does the paper report error bars suitably and correctly defined or other appropriate information about the statistical significance of the experiments?
    \item[] Answer: \answerYes{}
    \item[] Justification: We report Wilson 95\% confidence intervals on recovery rates (CA ours $n{=}200$/rule over 5 seeds; Sokoban $n{=}36$/size over 3 seeds), computed by the Wilson formula on binomial proportions; the small LLM sample sizes are stated explicitly.

\item {\bf Experiments compute resources}
    \item[] Question: For each experiment, does the paper provide sufficient information on the computer resources (type of compute workers, memory, time of execution) needed to reproduce the experiments?
    \item[] Answer: \answerYes{}
    \item[] Justification: The reasoner runs on CPU (a two-layer MLP plus exhaustive MAP enumeration, feasible to $L\approx 12$); the LLM baselines are API calls (o3 is minutes per call). Per-instance candidate counts and the enumeration wall are reported in Sections~\ref{sec:results} and~\ref{sec:limitations}.

\item {\bf Code of ethics}
    \item[] Question: Does the research conducted in the paper conform, in every respect, with the NeurIPS Code of Ethics \url{https://neurips.cc/public/EthicsGuidelines}?
    \item[] Answer: \answerYes{}
    \item[] Justification: The research conforms to the NeurIPS Code of Ethics; it uses only synthetic toy domains and public LLM APIs, with no human subjects or sensitive data.

\item {\bf Broader impacts}
    \item[] Question: Does the paper discuss both potential positive societal impacts and negative societal impacts of the work performed?
    \item[] Answer: \answerYes{}
    \item[] Justification: Section~\ref{sec:conclusion} notes the practical use (a missing evaluation axis, and a verifiable non-LLM inverter for diagnosis / root-cause / parameter-inference tasks). The work is foundational evaluation on toy domains with no direct path to harmful deployment.

\item {\bf Safeguards}
    \item[] Question: Does the paper describe safeguards that have been put in place for responsible release of data or models that have a high risk for misuse (e.g., pre-trained language models, image generators, or scraped datasets)?
    \item[] Answer: \answerNA{}
    \item[] Justification: The released assets are a small MLP and synthetic CA/Sokoban data, which pose no misuse risk.

\item {\bf Licenses for existing assets}
    \item[] Question: Are the creators or original owners of assets (e.g., code, data, models), used in the paper, properly credited and are the license and terms of use explicitly mentioned and properly respected?
    \item[] Answer: \answerYes{}
    \item[] Justification: We credit the LLMs evaluated (OpenAI o3 and GPT-4o-mini; Anthropic Claude Sonnet 4.6) and cite all prior work; the cellular automaton and Sokoban are standard, self-implemented domains.

\item {\bf New assets}
    \item[] Question: Are new assets introduced in the paper well documented and is the documentation provided alongside the assets?
    \item[] Answer: \answerYes{}
    \item[] Justification: The released code, instance generators, and cached transcripts are documented (README and \texttt{reproduce.sh}) under an MIT license.

\item {\bf Crowdsourcing and research with human subjects}
    \item[] Question: For crowdsourcing experiments and research with human subjects, does the paper include the full text of instructions given to participants and screenshots, if applicable, as well as details about compensation (if any)?
    \item[] Answer: \answerNA{}
    \item[] Justification: The paper involves no crowdsourcing and no research with human subjects.

\item {\bf Institutional review board (IRB) approvals or equivalent for research with human subjects}
    \item[] Question: Does the paper describe potential risks incurred by study participants, whether such risks were disclosed to the subjects, and whether Institutional Review Board (IRB) approvals (or an equivalent approval/review based on the requirements of your country or institution) were obtained?
    \item[] Answer: \answerNA{}
    \item[] Justification: The paper involves no research with human subjects.

\item {\bf Declaration of LLM usage}
    \item[] Question: Does the paper describe the usage of LLMs if it is an important, original, or non-standard component of the core methods in this research? Note that if the LLM is used only for writing, editing, or formatting purposes and does \emph{not} impact the core methodology, scientific rigor, or originality of the research, declaration is not required.
    \item[] Answer: \answerYes{}
    \item[] Justification: Frontier LLMs (o3, Claude Sonnet 4.6, GPT-4o-mini) are the evaluated baseline, central to the paper and fully described in Section~\ref{sec:setup} and Appendix~\ref{app:prompts}; our reasoner itself is non-transformer and uses no LLM.

\end{enumerate}
